\section[Continuous Random Variables]{Chapter 5: Continuous Random Variables}

% ============================================================
\subsection{Probability Density Functions}
% ============================================================

\bigskip

A discrete r.v.\ takes values you can list. A \textbf{continuous} r.v.\ can take any value in an interval, like $(0, 1)$, $(0, \infty)$ or all of $\mathbb{R}$. Its CDF has no jumps, so instead of a PMF we describe it with a \emph{density}.

\bigskip

\begin{definition}[Continuous r.v.]\label{def:contrv}
An r.v.\ has a \textbf{continuous distribution} if its CDF is differentiable. We also allow finitely many points (like endpoints) where the CDF is continuous but not differentiable, as long as it's differentiable everywhere else. A \textbf{continuous r.v.} is an r.v.\ with a continuous distribution.
\end{definition}

\bigskip

\begin{definition}[Probability Density Function]\label{def:pdf}
For a continuous r.v.\ $X$ with CDF $F$, the \textbf{probability density function} (PDF) of $X$ is the derivative of the CDF:
\medskip
\[
f(x) = F'(x).
\]
\medskip
The \textbf{support} of $X$ is the set of all $x$ where $f(x) > 0$.
\end{definition}

\bigskip

\begin{remark}
For a continuous r.v., $P(X = x) = 0$ for \emph{every} $x$. The probability of a single value is the height of the jump in the CDF at that value, and a continuous CDF has no jumps. That's why a PMF would be useless here (it would be $0$ everywhere) and we use a PDF instead.
\end{remark}

\bigskip

\begin{intuition}
$f(x)$ is \emph{not} a probability, and it can be bigger than $1$. It's a probability \emph{per unit length}. For a tiny interval of length $\varepsilon$ around $x$,
\medskip
\[
P\!\left(x - \tfrac{\varepsilon}{2} < X < x + \tfrac{\varepsilon}{2}\right) = \int_{x - \varepsilon/2}^{x + \varepsilon/2} f(t)\,dt \approx f(x)\,\varepsilon,
\]
\medskip
since $f$ is roughly constant on a tiny interval. So think of $f(x)\,dx$ as ``the probability that $X$ lands in a tiny interval of length $dx$ at $x$''. Units check: if $X$ is in cm, then $f(x)$ is in cm$^{-1}$, and multiplying by $dx$ (in cm) gives a unitless probability.
\end{intuition}

\bigskip

\begin{proposition}[PDF to CDF]\label{prop:pdf2cdf}
Let $X$ be a continuous r.v.\ with PDF $f$. Then its CDF is
\medskip
\[
F(x) = \int_{-\infty}^{x} f(t)\,dt.
\]
\end{proposition}

\medskip

\begin{proof}
By definition, $F$ is an antiderivative of $f$, so by the fundamental theorem of calculus,
\medskip
\begin{align*}
\int_{-\infty}^{x} f(t)\,dt
  &= F(x) - \lim_{t \to -\infty} F(t)
  && \by{FTC} \\[0.6em]
  &= F(x) - 0 = F(x).
  && \by{CDF $\to 0$ at $-\infty$}
\end{align*}
\end{proof}

\bigskip

\begin{intuition}
This mirrors the discrete case, where $F(x)$ was the sum of the PMF over all values $\le x$. Here $F(x)$ is the \emph{area} under the PDF to the left of $x$. Differentiating and integrating undo each other, so the PDF and CDF carry exactly the same information.
\end{intuition}

\bigskip

\begin{formula}[Getting probabilities from a PDF]
To get a probability, \textbf{integrate the PDF over the region you care about}:
\medskip
\begin{align*}
P(a < X \le b) &= F(b) - F(a) = \int_{a}^{b} f(x)\,dx,
  && \by{interval} \\[0.6em]
P(X \in A) &= \int_{A} f(x)\,dx.
  && \by{any region $A \subseteq \mathbb{R}$}
\end{align*}
\medskip
Since single points have probability $0$, endpoints don't matter:
\medskip
\[
P(a < X < b) = P(a < X \le b) = P(a \le X < b) = P(a \le X \le b).
\]
\end{formula}

\bigskip

\begin{pitfall}
Being carefree about endpoints is only OK for \emph{continuous} r.v.s. For a discrete r.v., $P(X \le 3)$ and $P(X < 3)$ can be different, since $P(X = 3)$ can be positive.
\end{pitfall}

\bigskip

\begin{theorem}[Valid PDFs]\label{thm:validpdf}
The PDF $f$ of a continuous r.v.\ must satisfy:
\medskip
\begin{enumerate}
  \setlength{\itemsep}{0.4em}
  \item \textbf{Nonnegative:} $f(x) \ge 0$ for all $x$;
  \item \textbf{Integrates to 1:} $\displaystyle\int_{-\infty}^{\infty} f(x)\,dx = 1$.
\end{enumerate}
\medskip
Conversely, any function with these two properties is the PDF of some continuous r.v.
\end{theorem}

\medskip

\begin{proof}
(1) $f(x)$ is the slope of the CDF at $x$, and a CDF can never decrease, so its slope can't be negative. (Or: if $f(x_0) < 0$, integrating over a tiny interval around $x_0$ would give a negative probability.) (2) The integral over all of $\mathbb{R}$ is $P(-\infty < X < \infty) = 1$.
\end{proof}

\bigskip

\begin{example}[Finding the constant, and a density above 1]\label{ex:cx2}
Let $f(x) = c\,x^{2}$ for $0 \le x \le 2$, and $f(x) = 0$ otherwise. Find $c$ so that $f$ is a valid PDF, then find the CDF and $P(X > 1)$.
\end{example}

\medskip

\begin{solution}
For $c > 0$, $f$ is nonnegative, so we just need the total area to be $1$:
\medskip
\begin{align*}
1 &= \int_{0}^{2} c\,x^{2}\,dx
  && \by{integrates to 1} \\[0.6em]
  &= c \cdot \frac{x^{3}}{3}\bigg|_{0}^{2} = \frac{8c}{3},
\end{align*}
\medskip
so $c = \tfrac{3}{8}$. Notice $f(2) = \tfrac{3}{8} \cdot 4 = \tfrac{3}{2} > 1$, which is fine for a density.

\medskip

For $0 \le x \le 2$, the CDF is
\medskip
\[
F(x) = \int_{0}^{x} \frac{3}{8} t^{2}\,dt = \frac{x^{3}}{8},
\]
\medskip
with $F(x) = 0$ for $x < 0$ and $F(x) = 1$ for $x > 2$. Then
\medskip
\[
P(X > 1) = 1 - F(1) = 1 - \frac{1}{8} = \frac{7}{8}.
\]
\end{solution}

\bigskip

\begin{example}[Logistic distribution]\label{ex:logistic}
The \textbf{Logistic} distribution has CDF
\medskip
\[
F(x) = \frac{e^{x}}{1 + e^{x}}, \qquad x \in \mathbb{R}.
\]
\medskip
Find the PDF, and find $P(-2 < X < 2)$ for $X \sim$ Logistic.
\end{example}

\medskip

\begin{solution}
Differentiate the CDF with the quotient rule:
\medskip
\[
f(x) = \frac{e^{x}(1 + e^{x}) - e^{x} \cdot e^{x}}{(1 + e^{x})^{2}} = \frac{e^{x}}{(1 + e^{x})^{2}}, \qquad x \in \mathbb{R}.
\]
\medskip
\begin{center}
\begin{tikzpicture}[x=0.8cm, y=12cm, >=stealth]
  \fill[gray!25, domain=-2:2, samples=60]
    (-2,0) -- plot (\x, {1/((exp(\x/2)+exp(-\x/2))^2)}) -- (2,0) -- cycle;
  \draw[->, gray!70!black] (-5.3,0) -- (5.5,0) node[right, black] {$x$};
  \draw[->, gray!70!black] (0,0) -- (0,0.31) node[above, black] {$f(x)$};
  \draw[ink, thick, domain=-5:5, samples=120]
    plot (\x, {1/((exp(\x/2)+exp(-\x/2))^2)});
  \foreach \x in {-4,-2,2,4}
    \draw[gray!70!black] (\x,0.006) -- (\x,-0.006) node[below, black] {$\x$};
  \draw[gray, dashed] (-2.6,0.25) -- (0,0.25);
  \node[left] at (-2.6,0.25) {$\tfrac14$};
  \node[ink, fill=gray!25, inner sep=1.5pt] at (0,0.07) {\small $\approx 0.76$};
\end{tikzpicture}
\end{center}
\medskip
The shaded area is $P(-2 < X < 2)$. Since we already have $F$ as an antiderivative,
\medskip
\[
P(-2 < X < 2) = F(2) - F(-2) = \frac{e^{2}}{1 + e^{2}} - \frac{e^{-2}}{1 + e^{-2}} \approx 0.881 - 0.119 = 0.76.
\]
\bigskip

\textbf{Without knowing $F$.} Substitute $u = 1 + e^{x}$, so $du = e^{x}\,dx$, and the limits $x = \pm 2$ become $u = 1 + e^{\pm 2}$:
\medskip
\[
\int_{-2}^{2} \frac{e^{x}}{(1 + e^{x})^{2}}\,dx
= \int_{1 + e^{-2}}^{1 + e^{2}} \frac{1}{u^{2}}\,du
= \left[-\frac{1}{u}\right]_{1 + e^{-2}}^{1 + e^{2}}
\approx 0.76.
\]
\end{solution}

\bigskip

\begin{example}[Rayleigh distribution]\label{ex:rayleigh}
The \textbf{Rayleigh} distribution has CDF
\medskip
\[
F(x) = 1 - e^{-x^{2}/2}, \qquad x > 0.
\]
\medskip
Find the PDF and $P(X > 2)$.
\end{example}

\medskip

\begin{solution}
By the chain rule,
\medskip
\[
f(x) = F'(x) = -e^{-x^{2}/2} \cdot (-x) = x\,e^{-x^{2}/2}, \qquad x > 0.
\]
\medskip
We could integrate $f$ from $2$ to $\infty$ with $u = x^{2}/2$, but $F$ already does the work:
\medskip
\[
P(X > 2) = \int_{2}^{\infty} x\,e^{-x^{2}/2}\,dx = 1 - F(2) = e^{-2} \approx 0.135.
\]
\end{solution}

\bigskip

\begin{definition}[Expectation of a Continuous r.v.]\label{def:contE}
The \textbf{expected value} of a continuous r.v.\ $X$ with PDF $f$ is
\medskip
\[
\E(X) = \int_{-\infty}^{\infty} x\,f(x)\,dx,
\]
\medskip
as long as $\int_{-\infty}^{\infty} |x|\,f(x)\,dx < \infty$. In practice we only integrate over the support, since $f = 0$ elsewhere.
\end{definition}

\bigskip

\begin{intuition}
Same idea as the discrete case: replace the sum with an integral and the PMF with the PDF. $\E(X)$ is still the balancing point of the distribution. The units work out too: $x\,f(x)\,dx$ has units cm $\cdot$ cm$^{-1}$ $\cdot$ cm $=$ cm.
\end{intuition}

\bigskip

\begin{example}[Mean of the $x^{2}$ density]\label{ex:cx2mean}
Find $\E(X)$ for the PDF $f(x) = \tfrac{3}{8}x^{2}$ on $[0, 2]$ from Example~\ref{ex:cx2}.
\end{example}

\medskip

\begin{solution}
\medskip
\[
\E(X) = \int_{0}^{2} x \cdot \frac{3}{8}x^{2}\,dx = \frac{3}{8} \cdot \frac{x^{4}}{4}\bigg|_{0}^{2} = \frac{3}{8} \cdot 4 = \frac{3}{2}.
\]
\medskip
The mean is above the midpoint $1$, which makes sense since the density piles up toward $2$.
\end{solution}

\bigskip

\begin{example}[Mean of the Logistic, by symmetry]
Show that $X \sim$ Logistic has $\E(X) = 0$.
\end{example}

\medskip

\begin{solution}
The PDF is symmetric about $0$: multiplying top and bottom by $e^{2x}$,
\medskip
\[
f(-x) = \frac{e^{-x}}{(1 + e^{-x})^{2}} = \frac{e^{x}}{(e^{x} + 1)^{2}} = f(x).
\]
\medskip
The tails of $f$ shrink like $e^{-|x|}$, so $\int |x|\,f(x)\,dx < \infty$ and the mean exists. For a symmetric density, the contributions $x\,f(x)$ and $(-x)\,f(-x)$ cancel, so $\E(X) = 0$.
\end{solution}

\bigskip

\begin{theorem}[LOTUS, Continuous]\label{thm:lotuscont}
If $X$ is a continuous r.v.\ with PDF $f$ and $g : \mathbb{R} \to \mathbb{R}$, then
\medskip
\[
\E\bigl(g(X)\bigr) = \int_{-\infty}^{\infty} g(x)\,f(x)\,dx.
\]
\end{theorem}

\medskip

\begin{proof}
We prove the case where $g$ is strictly increasing and differentiable. Let $Y = g(X)$, and let $(a, b)$ be the range of $g$ (possibly infinite). Since $g$ is strictly increasing, $g^{-1}$ exists and is strictly increasing too, so for $y \in (a, b)$,
\medskip
\begin{align*}
F_Y(y)
  &= P(g(X) \le y) = P\bigl(X \le g^{-1}(y)\bigr) = F_X\bigl(g^{-1}(y)\bigr).
  && \by{$g$ increasing}
\end{align*}
\medskip
Differentiating with the chain rule,
\medskip
\[
f_Y(y) = f_X\bigl(g^{-1}(y)\bigr) \cdot \frac{d}{dy}\, g^{-1}(y), \qquad a < y < b,
\]
\medskip
and $f_Y(y) = 0$ outside $(a, b)$. So
\medskip
\[
\E(Y) = \int_{a}^{b} y\, f_X\bigl(g^{-1}(y)\bigr) \frac{d}{dy}\, g^{-1}(y)\,dy.
\]
\medskip
Now substitute $x = g^{-1}(y)$. Then $y = g(x)$, $dx = \frac{d}{dy}g^{-1}(y)\,dy$, and as $y$ runs from $a$ to $b$, $x$ runs from $-\infty$ to $\infty$:
\medskip
\[
\E\bigl(g(X)\bigr) = \int_{-\infty}^{\infty} g(x)\,f_X(x)\,dx.
\]
\medskip
The general case (any $g$) holds too, but needs more machinery.
\end{proof}

\bigskip

\begin{remark}
Variance is defined exactly as before, $\Var(X) = \E\bigl[(X - \E X)^{2}\bigr]$, and the shortcut $\Var(X) = \E(X^{2}) - (\E X)^{2}$ still holds. With continuous LOTUS, $\E(X^{2}) = \int x^{2} f(x)\,dx$.
\end{remark}

\bigskip

\begin{example}[Variance of the $x^{2}$ density]
Find $\Var(X)$ for $f(x) = \tfrac{3}{8}x^{2}$ on $[0, 2]$.
\end{example}

\medskip

\begin{solution}
By continuous LOTUS with $g(x) = x^{2}$,
\medskip
\[
\E(X^{2}) = \int_{0}^{2} x^{2} \cdot \frac{3}{8}x^{2}\,dx = \frac{3}{8} \cdot \frac{x^{5}}{5}\bigg|_{0}^{2} = \frac{3}{8} \cdot \frac{32}{5} = \frac{12}{5}.
\]
\medskip
Using $\E(X) = \tfrac{3}{2}$ from Example~\ref{ex:cx2mean},
\medskip
\[
\Var(X) = \E(X^{2}) - (\E X)^{2} = \frac{12}{5} - \frac{9}{4} = \frac{48 - 45}{20} = \frac{3}{20}.
\]
\end{solution}

\medskip

% ============================================================
\subsection{Uniform}
% ============================================================

A uniform r.v.\ on the interval $(a, b)$ is a completely random number between $a$ and $b$.

\begin{definition}[Uniform Distribution]\label{def:unif}
  A continuous r.v.\ $U$ has the \textbf{Uniform distribution} on $(a, b)$, written $U \sim \Unif(a, b)$, if its PDF is
  \[
    f(x) =
    \begin{cases}
      \dfrac{1}{b - a}, & a < x < b, \\[0.6em]
      0,                & \text{otherwise}.
    \end{cases}
  \]
  The CDF is the accumulated area under the PDF:
  \[
    F(x) =
    \begin{cases}
      0,                  & x \le a,     \\[0.2em]
      \dfrac{x - a}{b - a}, & a < x < b, \\[0.6em]
      1,                  & x \ge b.
    \end{cases}
  \]
  $\Unif(0, 1)$ is called the \textbf{standard Uniform} distribution.
\end{definition}

\begin{center}
  \begin{tikzpicture}[x=1.6cm, y=1.6cm, >=stealth]
    % PDF
    \begin{scope}
      \draw[->, gray!70!black] (-0.6,0) -- (1.6,0) node[right, black] {$x$};
      \draw[->, gray!70!black] (0,0) -- (0,1.4) node[above, black] {PDF};
      \draw[ink, thick] (-0.5,0) -- (0,0);
      \draw[ink, thick] (0,1) -- (1,1);
      \draw[ink, thick] (1,0) -- (1.5,0);
      \draw[gray, dashed] (0,0) -- (0,1) (1,0) -- (1,1);
      \node[below] at (1,0) {$1$};
      \node[left] at (0,1) {$1$};
    \end{scope}

    % CDF
    \begin{scope}[xshift=4.6cm]
      \draw[->, gray!70!black] (-0.6,0) -- (1.6,0) node[right, black] {$x$};
      \draw[->, gray!70!black] (0,0) -- (0,1.4) node[above, black] {CDF};
      \draw[ink, thick] (-0.5,0) -- (0,0) -- (1,1) -- (1.5,1);
      \draw[gray, dashed] (1,0) -- (1,1);
      \node[below] at (1,0) {$1$};
      \node[left] at (0,1) {$1$};
    \end{scope}
  \end{tikzpicture}

  {\small $\Unif(0,1)$ PDF (left) and CDF (right).}
\end{center}

\begin{remark}
  For $U \sim \Unif(a, b)$ and any subinterval $(c, d)$ of $(a, b)$,
  \[
    P\bigl(U \in (c, d)\bigr) = \frac{d - c}{b - a},
  \]
  so probability is \emph{proportional to length}. For $\Unif(0, 1)$, probability \emph{is} length.
\end{remark}

\begin{proposition}[Conditioning a Uniform]\label{prop:unifcond}
  Let $U \sim \Unif(a, b)$ and let $(c, d)$ be a subinterval of $(a, b)$. Then the conditional distribution of $U$ given $U \in (c, d)$ is $\Unif(c, d)$.
\end{proposition}

\begin{proof}
  For $c < u < d$,
  \begin{align*}
    P\bigl(U \le u \mid U \in (c, d)\bigr)
      &= \frac{P\bigl(U \le u,\; c < U < d\bigr)}{P\bigl(U \in (c, d)\bigr)}
      && \by{def.\ of cond.\ prob.} \\[0.4em]
      &= \frac{P\bigl(U \in (c, u]\bigr)}{P\bigl(U \in (c, d)\bigr)}
      && \by{both events at once} \\[0.4em]
      &= \frac{(u - c)/(b - a)}{(d - c)/(b - a)}
       = \frac{u - c}{d - c},
      && \by{prob.\ $\propto$ length}
  \end{align*}
  which is the $\Unif(c, d)$ CDF.
\end{proof}

\begin{example}[Zooming in on a Uniform]
  Let $U \sim \Unif(0, 1)$. Find $P(0.2 < U < 0.6)$, and then $P(U < 0.5 \mid 0.2 < U < 0.6)$.
\end{example}

\begin{solution}
  Probability is length, so $P(0.2 < U < 0.6) = 0.4$. Given $U \in (0.2, 0.6)$, $U \sim \Unif(0.2, 0.6)$ by Proposition~\ref{prop:unifcond}, so
  \[
    P(U < 0.5 \mid 0.2 < U < 0.6) = \frac{0.5 - 0.2}{0.6 - 0.2} = \frac{0.3}{0.4} = \frac{3}{4}.
  \]
\end{solution}

\begin{theorem}[Uniform Mean and Variance]\label{thm:unifEV}
  If $U \sim \Unif(a, b)$, then
  \[
    \E(U) = \frac{a + b}{2},
    \qquad
    \Var(U) = \frac{(b - a)^{2}}{12}.
  \]
\end{theorem}

\begin{proof}
  The mean is the midpoint, as expected since the PDF is flat:
  \[
    \E(U) = \int_{a}^{b} x \cdot \frac{1}{b - a}\,dx = \frac{1}{b - a}\left(\frac{b^{2}}{2} - \frac{a^{2}}{2}\right) = \frac{a + b}{2}.
  \]
  For the variance, first find $\E(U^{2})$ with continuous LOTUS:
  \begin{align*}
    \E(U^{2})
      &= \int_{a}^{b} x^{2} \cdot \frac{1}{b - a}\,dx
       = \frac{1}{3} \cdot \frac{b^{3} - a^{3}}{b - a}
      && \by{LOTUS} \\[0.4em]
      &= \frac{a^{2} + ab + b^{2}}{3}.
      && \by{$b^{3} - a^{3} = (b - a)(a^{2} + ab + b^{2})$}
  \end{align*}
  Then
  \begin{align*}
    \Var(U)
      &= \frac{a^{2} + ab + b^{2}}{3} - \frac{(a + b)^{2}}{4}
       = \frac{4a^{2} + 4ab + 4b^{2} - 3a^{2} - 6ab - 3b^{2}}{12}
      && \by{common denominator} \\[0.4em]
      &= \frac{a^{2} - 2ab + b^{2}}{12}
       = \frac{(b - a)^{2}}{12}.
  \end{align*}
\end{proof}

\begin{definition}[Location-Scale Transformation]\label{def:locscale}
  Let $X$ be an r.v.\ and $Y = \sigma X + \mu$, where $\sigma$ and $\mu$ are constants with $\sigma > 0$. We say $Y$ is obtained from $X$ by a \textbf{location-scale transformation}: $\mu$ controls how the location is changed (shift) and $\sigma$ controls how the scale is changed (stretch).
\end{definition}

\begin{example}[Uniform mean and variance, the easy way]
  Use a location-scale transformation to get Theorem~\ref{thm:unifEV} without the algebra.
\end{example}

\begin{solution}
  Start with $U \sim \Unif(0, 1)$, whose PDF is $1$ on $(0, 1)$:
  \[
    \E(U) = \int_{0}^{1} x\,dx = \frac{1}{2},
    \qquad
    \E(U^{2}) = \int_{0}^{1} x^{2}\,dx = \frac{1}{3},
    \qquad
    \Var(U) = \frac{1}{3} - \frac{1}{4} = \frac{1}{12}.
  \]
  To turn $U$ into a general $\Unif(a, b)$ r.v., stretch the support from length $1$ to length $b - a$ (multiply by $b - a$, giving $\Unif(0, b - a)$), then shift it to start at $a$:
  \[
    \tilde{U} = a + (b - a)U \sim \Unif(a, b).
  \]
  Now linearity and the properties of variance do the work:
  \begin{align*}
    \E(\tilde{U}) &= a + (b - a)\,\E(U) = a + \frac{b - a}{2} = \frac{a + b}{2},
      && \by{linearity} \\[0.4em]
    \Var(\tilde{U}) &= (b - a)^{2}\,\Var(U) = \frac{(b - a)^{2}}{12}.
      && \by{shift drops out, scale comes out squared}
  \end{align*}
\end{solution}

% ============================================================
\subsection{Universality of the Uniform}
% ============================================================

Given a $\Unif(0, 1)$ r.v., we can construct an r.v.\ with any continuous distribution we want. Conversely, given an r.v.\ with any continuous distribution, we can create a $\Unif(0, 1)$ r.v. We call this the \textbf{universality of the Uniform}, because it tells us the Uniform is a universal starting point for building r.v.s with other distributions. (It's also called the probability integral transform, inverse transform sampling, or the quantile transformation.)

\begin{theorem}[Universality of the Uniform]\label{thm:universality}
  Let $F$ be a CDF that is continuous and strictly increasing on the support, so $F^{-1}$ exists, and let $U \sim \Unif(0, 1)$.
  \begin{enumerate}
    \item If $X = F^{-1}(U)$, then $X$ has CDF $F$.
    \item If $X$ has CDF $F$, then $F(X) \sim \Unif(0, 1)$.
  \end{enumerate}
\end{theorem}

\begin{proof}
  (1) For any $x$,
  \begin{align*}
    P(X \le x)
      &= P\bigl(F^{-1}(U) \le x\bigr)
       = P\bigl(U \le F(x)\bigr)
      && \by{apply $F$ to both sides; $F$ increasing} \\[0.4em]
      &= F(x).
      && \by{$P(U \le u) = u$ for $u \in (0, 1)$}
  \end{align*}
  (2) For $0 < u < 1$,
  \[
    P\bigl(F(X) \le u\bigr) = P\bigl(X \le F^{-1}(u)\bigr) = F\bigl(F^{-1}(u)\bigr) = u,
  \]
  which is the $\Unif(0, 1)$ CDF.
\end{proof}

\begin{intuition}
  Recipe for part (1): take the CDF you're interested in, take its inverse, and plug in the Uniform. $U$ picks a random height between $0$ and $1$ on the CDF's vertical axis, and $F^{-1}$ reads off which $x$ is at that height. Steep parts of $F$ (where the PDF is big) cover lots of heights over a short stretch of $x$, so $X$ lands there more often.
\end{intuition}

\begin{example}[Simulating a Logistic]
  Using one $U \sim \Unif(0, 1)$, construct an r.v.\ with the Logistic distribution from Example~\ref{ex:logistic}.
\end{example}

\begin{solution}
  Invert $F(x) = \frac{e^{x}}{1 + e^{x}}$ by solving $u = \frac{e^{x}}{1 + e^{x}}$ for $x$:
  \[
    u(1 + e^{x}) = e^{x}
    \;\Longrightarrow\;
    e^{x} = \frac{u}{1 - u}
    \;\Longrightarrow\;
    F^{-1}(u) = \ln\!\left(\frac{u}{1 - u}\right).
  \]
  So by Theorem~\ref{thm:universality}, $X = \ln\!\left(\frac{U}{1 - U}\right) \sim$ Logistic.
\end{solution}

% ============================================================
\subsection{Normal}
% ============================================================

The Normal is a continuous distribution with a bell-shaped PDF. Once we have the \emph{standard} Normal, we'll be able to get any Normal distribution we want by shifting and scaling.

\begin{definition}[Standard Normal Distribution]\label{def:stdnormal}
  A continuous r.v.\ $Z$ has the \textbf{standard Normal distribution}, written $Z \sim \Norm(0, 1)$, if its PDF is
  \[
    \varphi(z) = \frac{1}{\sqrt{2\pi}}\, e^{-z^{2}/2}, \qquad -\infty < z < \infty.
  \]
  The standard Normal CDF $\Phi$ is the accumulated area under the PDF:
  \[
    \Phi(z) = \int_{-\infty}^{z} \varphi(t)\,dt = \int_{-\infty}^{z} \frac{1}{\sqrt{2\pi}}\, e^{-t^{2}/2}\,dt.
  \]
\end{definition}

\begin{remark}
  There is no closed-form expression for the antiderivative of $\varphi$: we can't write $\Phi$ in terms of elementary functions. So $\Phi$ is just \emph{defined} as this integral, and in practice we get its values from software or a table.
\end{remark}

\begin{center}
  \begin{tikzpicture}[x=1cm, y=6cm, >=stealth]
    \fill[gray!30, domain=-1:1, samples=60]
      (-1,0) -- plot (\x, {exp(-\x*\x/2)/sqrt(2*pi)}) -- (1,0) -- cycle;
    \fill[gray!15, domain=-2:-1, samples=40]
      (-2,0) -- plot (\x, {exp(-\x*\x/2)/sqrt(2*pi)}) -- (-1,0) -- cycle;
    \fill[gray!15, domain=1:2, samples=40]
      (1,0) -- plot (\x, {exp(-\x*\x/2)/sqrt(2*pi)}) -- (2,0) -- cycle;
    \draw[->, gray!70!black] (-3.5,0) -- (3.7,0) node[right, black] {$z$};
    \draw[ink, thick, domain=-3.3:3.3, samples=120]
      plot (\x, {exp(-\x*\x/2)/sqrt(2*pi)});
    \foreach \x in {-3,-2,-1,0,1,2,3}
      \draw[gray!70!black] (\x,0.01) -- (\x,-0.01) node[below, black] {$\x$};
    \node[ink] at (0,0.15) {\small $\approx 68\%$};
    \node[ink] at (-1.5,0.035) {\scriptsize $13.5\%$};
    \node[ink] at (1.5,0.035) {\scriptsize $13.5\%$};
  \end{tikzpicture}

  {\small Standard Normal PDF $\varphi$. Dark: within $1$ SD of the mean; dark + light: within $2$ SDs ($\approx 95\%$).}
\end{center}

\begin{formula}[Symmetry properties of the standard Normal]
  \begin{enumerate}
    \item \textbf{Symmetry of the PDF:} $\varphi(z) = \varphi(-z)$.
    \item \textbf{Symmetry of tail areas:} the area to the left of $-z$ equals the area to the right of $z$, i.e.\ $P(Z \le -z) = P(Z \ge z)$, so
    \[
      \Phi(z) = 1 - \Phi(-z) \quad \text{for all } z.
    \]
    \item \textbf{Symmetry of $Z$ and $-Z$:} if $Z \sim \Norm(0, 1)$, then $-Z \sim \Norm(0, 1)$ as well.
  \end{enumerate}
\end{formula}

\begin{proof}
  (1) $z$ only appears as $z^{2}$. (2) Substitute $u = -t$ and use that the PDF integrates to $1$:
  \[
    \Phi(-z) = \int_{-\infty}^{-z} \varphi(t)\,dt = \int_{z}^{\infty} \varphi(u)\,du = 1 - \int_{-\infty}^{z} \varphi(u)\,du = 1 - \Phi(z).
  \]
  (3) $P(-Z \le z) = P(Z \ge -z) = 1 - \Phi(-z) = \Phi(z)$, using (2). So $-Z$ has CDF $\Phi$.
\end{proof}

\begin{theorem}[$\varphi$ is a Valid PDF]\label{thm:phivalid}
  $\displaystyle\int_{-\infty}^{\infty} e^{-z^{2}/2}\,dz = \sqrt{2\pi}$, so $\varphi$ integrates to $1$.
\end{theorem}

\begin{proof}
  We can't integrate $e^{-z^{2}/2}$ directly, so we square the integral and switch to polar coordinates. Since $z$ is a dummy variable, call it $x$ in one copy and $y$ in the other:
  \begin{align*}
    \left(\int_{-\infty}^{\infty} e^{-z^{2}/2}\,dz\right)^{2}
      &= \left(\int_{-\infty}^{\infty} e^{-x^{2}/2}\,dx\right)\left(\int_{-\infty}^{\infty} e^{-y^{2}/2}\,dy\right)
      && \by{rename dummy variables} \\[0.4em]
      &= \int_{-\infty}^{\infty}\int_{-\infty}^{\infty} e^{-(x^{2} + y^{2})/2}\,dx\,dy
      && \by{combine into a double integral} \\[0.4em]
      &= \int_{0}^{2\pi}\int_{0}^{\infty} e^{-r^{2}/2}\, r\,dr\,d\theta
      && \by{polar; extra $r$ is the Jacobian} \\[0.4em]
      &= \int_{0}^{2\pi}\left(\int_{0}^{\infty} e^{-u}\,du\right) d\theta
      && \by{$u = r^{2}/2$, $du = r\,dr$} \\[0.4em]
      &= \int_{0}^{2\pi} 1\,d\theta = 2\pi.
  \end{align*}
  Taking square roots, the integral is $\sqrt{2\pi}$.
\end{proof}

\begin{theorem}[Standard Normal Mean and Variance]\label{thm:stdnormalEV}
  If $Z \sim \Norm(0, 1)$, then $\E(Z) = 0$ and $\Var(Z) = 1$.
\end{theorem}

\begin{proof}
  \textbf{Mean.} $g(z) = z\,e^{-z^{2}/2}$ is an odd function ($g(-z) = -g(z)$), and the integral of an odd function over a symmetric range is $0$:
  \[
    \E(Z) = \frac{1}{\sqrt{2\pi}} \int_{-\infty}^{\infty} z\, e^{-z^{2}/2}\,dz = 0.
  \]
  \textbf{Variance.} Since $\E(Z) = 0$, $\Var(Z) = \E(Z^{2})$. The integrand $z^{2}e^{-z^{2}/2}$ is even, so
  \begin{align*}
    \Var(Z)
      &= \frac{1}{\sqrt{2\pi}} \int_{-\infty}^{\infty} z^{2} e^{-z^{2}/2}\,dz
       = \frac{2}{\sqrt{2\pi}} \int_{0}^{\infty} z^{2} e^{-z^{2}/2}\,dz
      && \by{LOTUS; even function} \\[0.4em]
      &= \frac{2}{\sqrt{2\pi}} \left( \Bigl[-z\,e^{-z^{2}/2}\Bigr]_{0}^{\infty} + \int_{0}^{\infty} e^{-z^{2}/2}\,dz \right)
      && \by{parts: $u = z$, $dv = z e^{-z^{2}/2}dz$} \\[0.4em]
      &= \frac{2}{\sqrt{2\pi}} \left( 0 + \frac{\sqrt{2\pi}}{2} \right) = 1.
      && \by{Thm~\ref{thm:phivalid}, half the area}
  \end{align*}
\end{proof}

\begin{definition}[Normal Distribution]\label{def:normal}
  If $Z \sim \Norm(0, 1)$, then
  \[
    X = \mu + \sigma Z
  \]
  has the \textbf{Normal distribution} with mean $\mu$ and variance $\sigma^{2}$, for any real $\mu$ and $\sigma > 0$. We write $X \sim \Norm(\mu, \sigma^{2})$.
\end{definition}

\begin{remark}
  This is a location-scale transformation of $Z$, so by linearity and the variance properties, $\E(X) = \mu + \sigma \cdot 0 = \mu$ and $\Var(X) = \sigma^{2} \cdot 1 = \sigma^{2}$. Going the other way, \textbf{standardization}: if $X \sim \Norm(\mu, \sigma^{2})$ then
  \[
    Z = \frac{X - \mu}{\sigma} \sim \Norm(0, 1).
  \]
\end{remark}

\begin{theorem}[Normal CDF and PDF]\label{thm:normalcdf}
  Let $X \sim \Norm(\mu, \sigma^{2})$. Then the CDF and PDF of $X$ are
  \[
    F(x) = \Phi\!\left(\frac{x - \mu}{\sigma}\right),
    \qquad
    f(x) = \varphi\!\left(\frac{x - \mu}{\sigma}\right) \frac{1}{\sigma} = \frac{1}{\sqrt{2\pi}\,\sigma} \exp\!\left(-\frac{(x - \mu)^{2}}{2\sigma^{2}}\right).
  \]
\end{theorem}

\begin{proof}
  Standardize inside the probability:
  \begin{align*}
    F(x) = P(X \le x)
      &= P\!\left(\frac{X - \mu}{\sigma} \le \frac{x - \mu}{\sigma}\right)
      && \by{same event, $\sigma > 0$} \\[0.4em]
      &= \Phi\!\left(\frac{x - \mu}{\sigma}\right).
      && \by{standardized $X$ is $\Norm(0,1)$}
  \end{align*}
  Differentiating with the chain rule gives $f(x) = \varphi\!\left(\frac{x - \mu}{\sigma}\right)\cdot\frac{1}{\sigma}$, and plugging in the formula for $\varphi$ gives the last form.
\end{proof}

\begin{formula}[68--95--99.7\% rule]
  If $X \sim \Norm(\mu, \sigma^{2})$, then
  \[
    P(|X - \mu| < \sigma) \approx 0.68,
    \qquad
    P(|X - \mu| < 2\sigma) \approx 0.95,
    \qquad
    P(|X - \mu| < 3\sigma) \approx 0.997.
  \]
  That is, the chance of falling within $1$, $2$, $3$ SDs of the mean.
\end{formula}

\begin{example}[Using the 68--95--99.7 rule]\label{ex:normalrule}
  Let $X \sim \Norm(-1, 4)$. Approximate $P(|X| < 3)$.
\end{example}

\begin{solution}
  Here $\mu = -1$ and $\sigma = 2$. $|X| < 3$ is the same as $-3 < X < 3$. Standardize with $Z = \frac{X - (-1)}{2}$:
  \begin{align*}
    P(-3 < X < 3)
      &= P\!\left(\frac{-3 - (-1)}{2} < \frac{X - (-1)}{2} < \frac{3 - (-1)}{2}\right)
      && \by{standardize} \\[0.4em]
      &= P(-1 < Z < 2).
  \end{align*}
  The rule says $P(-1 < Z < 1) \approx 0.68$ and $P(-2 < Z < 2) \approx 0.95$, so the two bands between $1$ and $2$ SDs hold $0.95 - 0.68 = 0.27$ together, and by symmetry $P(1 < Z < 2) \approx 0.27/2$. So
  \[
    P(-1 < Z < 2) = P(-1 < Z < 1) + P(1 < Z < 2) \approx 0.68 + 0.135 = 0.815.
  \]
  The exact value is $\Phi(2) - \Phi(-1) \approx 0.8186$, so the rule does well.
\end{solution}

\begin{example}[Folded Normal]\label{ex:foldednormal}
  Let $Y = |Z|$ with $Z \sim \Norm(0, 1)$. Find (a) $\E(Y)$, (b) $\Var(Y)$, (c) the CDF and PDF of $Y$.
\end{example}

\begin{solution}
  (a) By LOTUS, and since $|z|\varphi(z)$ is even,
  \begin{align*}
    \E(Y)
      &= \int_{-\infty}^{\infty} |z|\, \frac{1}{\sqrt{2\pi}} e^{-z^{2}/2}\,dz
       = \frac{2}{\sqrt{2\pi}} \int_{0}^{\infty} z\, e^{-z^{2}/2}\,dz
      && \by{even function} \\[0.4em]
      &= \frac{2}{\sqrt{2\pi}} \Bigl[-e^{-z^{2}/2}\Bigr]_{0}^{\infty}
       = \frac{2}{\sqrt{2\pi}}
       = \sqrt{\frac{2}{\pi}} \approx 0.80.
  \end{align*}
  (b) $\E(Y^{2}) = \E(Z^{2}) = \Var(Z) = 1$, so
  \[
    \Var(Y) = \E(Y^{2}) - (\E Y)^{2} = 1 - \frac{2}{\pi} \approx 0.36.
  \]
  (c) For $y \ge 0$,
  \[
    F_Y(y) = P(|Z| \le y) = P(-y \le Z \le y) = \Phi(y) - \Phi(-y) = 2\Phi(y) - 1,
  \]
  using $\Phi(-y) = 1 - \Phi(y)$, and $F_Y(y) = 0$ for $y < 0$. Differentiating,
  \[
    f_Y(y) = 2\varphi(y) = \sqrt{\frac{2}{\pi}}\, e^{-y^{2}/2}, \qquad y > 0.
  \]
\end{solution}

% ============================================================
\subsection{Exponential}
% ============================================================

The Exponential is the continuous counterpart of the Geometric. A Geometric r.v.\ counts the number of failures before the first success in a sequence of Bernoulli trials. With the Exponential, we're waiting for a success in \emph{continuous time}, where successes arrive at a rate of $\lambda$ per unit of time. The average number of successes in a time interval of length $t$ is $\lambda t$, though the actual number varies randomly. An Exponential r.v.\ represents the waiting time until the first arrival.

\begin{definition}[Exponential Distribution]\label{def:expo}
  A continuous r.v.\ $X$ has the \textbf{Exponential distribution} with parameter (rate) $\lambda > 0$, written $X \sim \Expo(\lambda)$, if its PDF is
  \[
    f(x) = \lambda e^{-\lambda x}, \qquad x > 0.
  \]
  This is valid since $\int_{0}^{\infty} \lambda e^{-\lambda x}\,dx = 1$. The CDF is
  \[
    F(x) = \int_{0}^{x} \lambda e^{-\lambda t}\,dt = 1 - e^{-\lambda x}, \qquad x > 0.
  \]
\end{definition}

\begin{proposition}[Scaling to $\Expo(1)$]\label{prop:exposcale}
  If $X \sim \Expo(\lambda)$, then $Y = \lambda X \sim \Expo(1)$.
\end{proposition}

\begin{proof}
  For $y > 0$,
  \[
    P(Y \le y) = P\!\left(X \le \frac{y}{\lambda}\right) = 1 - e^{-\lambda \cdot y/\lambda} = 1 - e^{-y},
  \]
  which is the $\Expo(1)$ CDF.
\end{proof}

\begin{theorem}[Exponential Mean and Variance]\label{thm:expoEV}
  If $X \sim \Expo(\lambda)$, then
  \[
    \E(X) = \frac{1}{\lambda},
    \qquad
    \Var(X) = \frac{1}{\lambda^{2}}.
  \]
\end{theorem}

\begin{proof}
  First do $Y \sim \Expo(1)$, using integration by parts each time.

  \textbf{First moment.} With $u = y$, $dv = e^{-y}dy$ (so $du = dy$, $v = -e^{-y}$):
  \[
    \E(Y) = \int_{0}^{\infty} y\, e^{-y}\,dy = \Bigl[-y\,e^{-y}\Bigr]_{0}^{\infty} + \int_{0}^{\infty} e^{-y}\,dy = 0 + 1 = 1.
  \]
  \textbf{Second moment.} With $u = y^{2}$, $dv = e^{-y}dy$:
  \[
    \E(Y^{2}) = \int_{0}^{\infty} y^{2} e^{-y}\,dy = \Bigl[-y^{2} e^{-y}\Bigr]_{0}^{\infty} + 2\int_{0}^{\infty} y\, e^{-y}\,dy = 0 + 2\,\E(Y) = 2,
  \]
  so $\Var(Y) = 2 - 1^{2} = 1$.

  \textbf{General $\lambda$.} By Proposition~\ref{prop:exposcale}, $X = Y/\lambda$ with $Y \sim \Expo(1)$, so
  \[
    \E(X) = \frac{\E(Y)}{\lambda} = \frac{1}{\lambda},
    \qquad
    \Var(X) = \frac{\Var(Y)}{\lambda^{2}} = \frac{1}{\lambda^{2}}.
  \]
\end{proof}

\begin{pitfall}
  $\lambda$ is a \emph{rate}, not a mean. If buses come on average every $10$ minutes, then $\lambda = 1/10$ per minute, not $10$. (Some books parametrize the Exponential by its mean instead, so always check.)
\end{pitfall}

\begin{definition}[Memoryless Property]\label{def:memoryless}
  A continuous distribution on $(0, \infty)$ is \textbf{memoryless} if, for $X$ with that distribution,
  \[
    P(X \ge s + t \mid X \ge s) = P(X \ge t) \quad \text{for all } s, t \ge 0.
  \]
  Think of $X$ as a waiting time (e.g.\ for a phone call): if you've already waited $s$ minutes, the chance of waiting at least $t$ \emph{more} minutes is the same as if you'd just started.
\end{definition}

\begin{theorem}[The Exponential is Memoryless]\label{thm:expomemoryless}
  If $X \sim \Expo(\lambda)$, then $X$ has the memoryless property.
\end{theorem}

\begin{proof}
  Here $P(X \ge s) = 1 - P(X \le s) = e^{-\lambda s}$. Then
  \begin{align*}
    P(X \ge s + t \mid X \ge s)
      &= \frac{P(X \ge s + t,\; X \ge s)}{P(X \ge s)}
      && \by{def.\ of cond.\ prob.} \\[0.4em]
      &= \frac{P(X \ge s + t)}{P(X \ge s)}
      && \by{$X \ge s$ is redundant} \\[0.4em]
      &= \frac{e^{-\lambda(s + t)}}{e^{-\lambda s}} = e^{-\lambda t} = P(X \ge t).
  \end{align*}
\end{proof}

\begin{example}[Waiting for a bus]
  Bus arrivals at a stop follow a rate of one every $10$ minutes on average, so your wait is $X \sim \Expo(1/10)$ in minutes. You've already waited $10$ minutes. What's the probability you wait more than $15$ minutes in total?
\end{example}

\begin{solution}
  By memorylessness, having waited $10$ minutes changes nothing:
  \[
    P(X > 15 \mid X > 10) = P(X > 5) = e^{-5/10} = e^{-0.5} \approx 0.61.
  \]
  Your expected \emph{remaining} wait is still $1/\lambda = 10$ minutes, no matter how long you've been standing there.
\end{solution}

\begin{theorem}[Memoryless Implies Exponential]\label{thm:memorylessexpo}
  If $X$ is a positive continuous r.v.\ with the memoryless property, then $X \sim \Expo(\lambda)$ for some $\lambda > 0$.
\end{theorem}

\begin{proof}
  Let $F$ be the CDF of $X$ and $G(x) = P(X > x) = 1 - F(x)$. The memoryless property says
  \[
    G(s + t) = G(s)\,G(t) \quad \text{for all } s, t \ge 0.
  \]
  We solve for $G$.
  \begin{itemize}
    \item Taking $s = t$ gives $G(2t) = G(t)^{2}$, then $G(3t) = G(t)^{3}$, and in general $G(kt) = G(t)^{k}$ for positive integers $k$.
    \item Replacing $t$ by $t/k$: $G(t) = G(t/k)^{k}$, so $G(t/k) = G(t)^{1/k}$.
    \item Combining the two: $G\!\left(\frac{m}{n}t\right) = G(t)^{m/n}$ for all positive integers $m, n$, so $G(xt) = G(t)^{x}$ for every positive rational $x$.
    \item Any positive real number is a limit of positive rationals, and $G$ is continuous, so $G(xt) = G(t)^{x}$ for all real $x > 0$.
  \end{itemize}
  Taking $t = 1$: $G(x) = G(1)^{x} = e^{-\lambda x}$, where $\lambda = -\ln G(1)$. Since $X$ is positive and continuous, $0 < G(1) < 1$, so $\lambda > 0$. Then $F(x) = 1 - e^{-\lambda x}$, the $\Expo(\lambda)$ CDF.
\end{proof}


